\section{Monitoring and controlling reasoning}
\label{sec:p19-annex-rationality-8-6}

\begin{annexitem}{annex:p19.r86.stopping}{thm:p19-r86-stopping}{Capped threshold rules and increasing trajectories}
\begin{proof}
The finite index set contains the cap index \(m\). Taking the smallest
index in \(\{j:q_j\geq h_j\}\cup\{m\}\) therefore gives one
well-defined stopping index. If \(h_j'\leq h_j\) everywhere, every
threshold crossing for \(h\) is also a crossing for \(h'\); including
the same cap preserves this inclusion. The minimum for \(h'\) is thus
no larger. There is no correctness variable in that comparison, so assigning
either correctness outcome to the selected answer preserves the stopping
calculation. Threshold behavior alone cannot determine correctness.

For the continuous example,
\(.6+.04t\geq.9-.005t\) exactly when
\(.045t\geq.3\), whose first solution is \(20/3\).
At that time confidence is \(.6+.04(20/3)=13/15\).
The second inequality is \(.3+.005t\geq.9-.005t\), first satisfied
at \(t=60\), beyond the cap. It therefore stops at ten with confidence
\(.3+.005(10)=.35\). Both trajectories have positive slopes and stay
in \([0,1]\) on the stated interval. Across these two items more time
is associated with lower final confidence, without any within-item decline.
This is a construction, not an estimated psychological relation.
\end{proof}
\end{annexitem}

\begin{annexitem}{annex:p19.r86.trace-calculations}{thm:p19-r86-trace-calculations}{Complete calculations on the fictional record}
\begin{proof}
The ball equation gives \(x=.05\). For both entailment failures use the
valuation \(A=0,B=1\): the premises are true and the respective
conclusions \(A\) and \(\neg B\) are false. Union probability cannot
be below either constituent probability, so .5 is impossible. Every Boolean
valuation makes \(P\lor\neg P\) true. Of consecutive integers
\(m,m-1\), one is even, proving the final key. These establish all six
criterion outcomes without using confidence or duration.

For every item \(y_i^{(1)}-y_i^{(0)}\) equals one for a corrected
answer, minus one for a harmful change, and zero otherwise. Summing proves
the accuracy identity. The example has one correction (item 3) and one
harmful change (item 4). Both sums of correctness values are three.
Initial confidence sums to \(77/20\); final confidence sums to
\(24/5\). Division by six gives \(77/120\) and \(4/5\).
The extra times sum to \(20+40+30+35+5+60=190\).

The six initial squared losses are
\(.81,.2025,.25,.16,.04,.25\), summing to \(1.7125=137/80\).
Dividing by six gives \(137/480\). The final losses are
\(.9025,.0225,.09,.64,.01,.36\), summing to \(2.025=81/40\).
Their mean is \(27/80\), larger than the initial mean.

For initial gamma, the correct-answer judgments are .55, .60 and .80;
the incorrect-answer judgments are .90, .50 and .50. Each correct judgment
has two concordant and one discordant cross-outcome pairs, so
\(N_+=6,N_-=3\), giving gamma \(1/3\). Tied outcomes are excluded;
the equal .50 judgments occur within one outcome class and add no comparable
pair. For final gamma, correct judgments .85, .70 and .90 have, against
incorrect judgments .95, .80 and .60, respective concordant/discordant
counts \((2,1),(1,2),(2,1)\). Thus \(N_+=5,N_-=4\), giving
\(1/9\). These computations use the exact comparable-pair definition.

Final confidence at least .8 selects items 1, 2, 4 and 5. Two of them are
correct, so selected accuracy is \(2/4=1/2\), and coverage is
\(4/6=2/3\). Correct item 3 is withheld; incorrect item 6 is withheld
as well. The rule therefore has both commission and omission consequences
on this particular record. Its mean selected confidence is .875, which does
not equal its observed selected accuracy .5. No independent-trial sampling
or population law was assumed in these finite calculations.
\end{proof}
\end{annexitem}

\begin{annexitem}{annex:p19.r86.calibrated-control}{thm:p19-r86-calibrated-control}{Calibration and a reporting threshold}
\begin{proof}
The positive-mass event \(A=\{C\geq h\}\) is a union of confidence
classes. In the finite law,
\[
 \mathbb P(Y=1\mid A)
  ={\sum_{c\geq h}\mathbb P(C=c)\mathbb P(Y=1\mid C=c)
      \over\mathbb P(A)}
  ={\sum_{c\geq h}\mathbb P(C=c)c\over\mathbb P(A)}.
\]
The second equality uses calibration on each positive-mass class; null
classes contribute zero. This is \(E(C\mid A)\), at least \(h\)
because every included confidence is at least \(h\).

For a counterexample to a gamma-only guarantee, use the stipulated outcomes
\((0,0,1,1)\) and reports \((.99,.99,1,1)\). All four cross-outcome
pairs are concordant, giving gamma one as in
\ref{thm:p19-jol-gamma}. A threshold .95 nevertheless selects every
answer, with accuracy \(1/2<.95\). The zero-outcome .99 class is not
calibrated, so this does not contradict the theorem. It also shows why an
ordering property cannot replace a probability calibration premise. The
equation gives a conditional accuracy property, but contains no utilities
for omission or error and therefore proves no optimal threshold.
\end{proof}
\end{annexitem}

\begin{annexitem}{annex:p19.r86.step-value}{thm:p19-r86-step-value}{Expected gains under the declared belief and value model}
\begin{proof}
The current answer has expected value \(Rp_0-L(1-p_0)\).
The additional procedure has expected net value
\(Rp_1-L(1-p_1)-\kappa\). Their difference is
\((R+L)(p_1-p_0)-\kappa\). Its sign gives the two strict
preferences and its zero value gives indifference.

At \(R=1,L=0,p_0=.6\), improvement to .9 at cost .1 gives
\(.9-.6-.1=.2>0\); improvement to .61 gives
\(.61-.6-.1=-.09<0\). The same current probability thus supports
opposite actions under different prospective improvements. A deadline can
exclude an otherwise attractive procedure because it is infeasible; that
does not change the value identity for feasible options. The argument uses
the supplied belief model and utility/cost values. It does not prove their
calibration, describe actual human choices, or identify an optimal policy
for an unspecified future sequence of computations.
\end{proof}
\end{annexitem}
